\chapter{Deep Dive: Kafka Delivery Guarantees --- At-Least-Once, Idempotent Consumers, and Dead Letters}
\label{ch:dd-kafka}

Chapter~10 gave the rule of thumb --- Kafka is at-least-once, so make
consumers idempotent --- and Chapter~11 showed the poison-pill defense.
This deep dive reads the producer and consumer settings the project
actually runs, most of which it never wrote down, and follows one
message through every place it can be duplicated or lost. Then it adds
the two things a production consumer needs and the project lacks: a
deduplication record and a retry path that ends in a dead-letter topic
instead of a log line.

Three questions organize the chapter:

\begin{enumerate}
  \item \textbf{Where does the project sit} on the at-most-once /
  at-least-once / exactly-once scale --- not by intent, but by the
  defaults it inherited?
  \item \textbf{Exactly when is a message delivered twice}, and what
  must a consumer do to make that harmless?
  \item \textbf{What happens to a message that cannot be processed},
  and where should it go instead of nowhere?
\end{enumerate}

\section{Why this, and what it solves}

The project uses Kafka for one job: decoupling a business write from
its side effects. A user registers; a mail must go out; a student row
must appear in another database. Doing those inline would make
registration as slow and as fragile as SMTP and a second Postgres.
Doing them through a log makes them \emph{eventually} happen, survive
consumer restarts, and be replayable. Chapter~6 covered the sync-vs-async
split; what it did not say is that ``eventually'' has a precise
meaning in Kafka, made of five or six settings, and the defaults of
those settings are what the project is running today.

\begin{decision}[events, not commands --- and when a call is right after all]
\code{ts.user.registered} says ``this happened''; it does not say ``send
an e-mail''. Notification-service decides what a registration means
for it, and course-service ignores the topic entirely. That is the
event style, and it is why a fourth consumer can be added without
touching auth-service. The command style --- a \code{SendEmail} message
addressed to one handler --- couples the sender to the receiver's
existence and is usually a sign that a synchronous call was wanted. And
sometimes it was: when the caller needs the \emph{result} now (does this
user exist? is this file scanned?), a Feign call with a circuit breaker
is the honest design, and dressing it as request-reply over Kafka only
adds two hops and a correlation id. The project's rule: events for
facts other services react to; calls for questions the caller must have
answered before it can continue.
\end{decision}

\section{The settings the project never wrote}

\subsection{Producer}

Course-service builds its producer by hand in \code{KafkaProducerConfig}:

\begin{lstlisting}[language=Java]
Map<String, Object> configProps = new HashMap<>();
configProps.put(ProducerConfig.BOOTSTRAP_SERVERS_CONFIG, bootstrapServers);
configProps.put(ProducerConfig.KEY_SERIALIZER_CLASS_CONFIG,
        StringSerializer.class);                              (*@\co{1}@*)
configProps.put(ProducerConfig.VALUE_SERIALIZER_CLASS_CONFIG,
        JsonSerializer.class);
return new DefaultKafkaProducerFactory<>(configProps);
\end{lstlisting}

Three lines --- and everything else is the Kafka~3.9 client default.
Since Kafka~3.0 those defaults are the safe ones:

\begin{lstlisting}[language=properties]
acks=all
enable.idempotence=true
retries=2147483647
max.in.flight.requests.per.connection=5
delivery.timeout.ms=120000
max.block.ms=60000
linger.ms=0
\end{lstlisting}

Auth-service and notification-service use Spring Boot's auto-configured
factory from the YAML, which sets the same two serializers and nothing
else, so all three producers behave alike.

\begin{description}
  \item[\co{1}] The key serializer matters more than it looks. Every
  publisher in the project passes a key: \code{event.userId()} in
  \code{UserRegisteredPublisher} and \code{UserActivatedPublisher},
  \code{assignmentId} in \code{AssignmentCreatedPublisher}. Kafka hashes
  the key to choose a partition; equal keys, equal partition, and the
  broker preserves order within a partition. That is the entire
  per-user ordering guarantee, and it is one string.
\end{description}

What the defaults mean, one at a time:

\begin{itemize}
  \item \textbf{\code{acks=all}}: the leader waits for every in-sync
  replica before acknowledging. With \code{replicas: 1} in
  \code{kafka-topics.yaml} and a single broker in Compose, ``all'' is
  one --- durability equals one disk. In production this pairs with
  \code{replication.factor=3} and \code{min.insync.replicas=2}: a
  write survives one broker dying, and if two die the producer gets
  \code{NotEnoughReplicas} rather than a silent single-copy write.
  \item \textbf{Idempotent producer}: each producer gets an id and
  numbers its messages per partition; a retry after a lost ack is
  recognized by the broker and not appended twice. This closes the
  \emph{producer-side} duplicate from Chapter~10. It does nothing for
  the consumer side.
  \item \textbf{\code{retries} unbounded, bounded by
  \code{delivery.timeout.ms}}: the producer retries transient errors for
  up to two minutes, then fails the future. \code{send()} returned long
  ago; nobody in the project reads that future.
\end{itemize}

\begin{pitfall}[the YAML that does nothing]
Symptom: you add \code{spring.kafka.producer.acks: 1} to
\code{course-service.yml} to speed up a load test, and nothing changes.
Cause: \code{KafkaProducerConfig} declares its own
\code{ProducerFactory} bean, which switches off Boot's
auto-configuration for that factory --- the YAML producer block is read
by nobody in course-service. Auth-service, with no such bean,
\emph{would} honour it. Two services, one YAML convention, two
behaviours. Fix: delete the hand-written bean (Boot's factory does the
same thing) or move every property into it; never both.
\end{pitfall}

\subsection{Consumer}

The consumer YAML in \code{notification-service.yml} and
\code{course-service.yml} names a group, \code{auto-offset-reset:
earliest}, and the error-handling deserializers. Everything about
\emph{committing} is again a default, this time Spring Kafka's:

\begin{itemize}
  \item \code{enable.auto.commit=false}. Spring switches the client's
  own timer-based commit off and commits itself.
  \item \code{AckMode.BATCH}. After the listener has returned for
  \emph{every} record of one \code{poll()}, Spring commits the highest
  offset of the batch. Processing and commit are separated by the
  batch's processing time.
  \item \code{DefaultErrorHandler} with \code{FixedBackOff(0, 9)}: a
  listener that throws is retried nine more times with no delay, then
  the record is \emph{logged and skipped}. Not sent anywhere. Not
  counted anywhere the project reads.
  \item \code{max.poll.records=500}, \code{max.poll.interval.ms=300000}:
  a batch of up to 500 records must be fully processed within five
  minutes or the consumer is presumed dead.
\end{itemize}

Put those four together and the project's position is: at-least-once,
with duplicates possible on any crash or rebalance, and with silent loss
after ten failed attempts. Neither the ordering nor the loss is written
in any file you can \code{grep}.

\section{Exactly when a duplicate happens}

The timeline that produces one, with \code{notification-group} running
two instances A and B on a topic with two partitions:

\begin{lstlisting}
t0   A polls partition 0, gets offsets 40..49   (committed: 40)
t1   A processes 40..47: eight e-mails sent
t2   A is killed (rolling update, OOMKilled, SIGKILL)
     -- offset 40 is still the committed position --
t3   group rebalances; partition 0 is assigned to B
t4   B polls from the committed offset: 40
t5   B processes 40..49: ten e-mails sent
     -- users 40..47 receive their activation mail twice --
\end{lstlisting}

Nothing here is a bug. It is the contract: Kafka promises the offset
will not advance past a record you did not finish; the price is that a
record you \emph{did} finish, but did not commit, comes back. Every
consumer that does anything non-reversible --- send mail, insert a row,
charge a card --- must be written for \code{t5}.

Two variants matter for this project. \code{AckMode.RECORD} commits
after each record instead of each batch: fewer duplicates on a crash
(one instead of up to 500), more commit round-trips. And the
\emph{slow-consumer} variant needs no crash at all: if processing a
batch takes longer than \code{max.poll.interval.ms}, the broker removes
the consumer from the group, another instance takes the partition from
the last commit, and \emph{both} end up processing the same records ---
the first one has not noticed it was evicted yet.

\begin{pitfall}[one assignment, a storm of duplicate e-mails]
Symptom: a teacher creates an assignment in a course with 300 students;
some students receive the notification three times; the consumer log
shows \code{Member ... has left the group} and a rebalance. Cause:
\code{AssignmentCreatedListener} loops over
\code{enrolledStudentEmails} and calls \code{javaMailSender.send()} for
each, synchronously, on the poll thread. At one second per SMTP
round-trip, 300 e-mails is five minutes --- exactly
\code{max.poll.interval.ms}. The consumer is evicted mid-loop, the
partition is reassigned, and the new owner starts the loop again from
the first student. Fix, in order of cost: raise
\code{max.poll.interval.ms} (a patch), publish one
\code{ts.notification.email} record per recipient so each mail is its
own unit of work (the right shape), and make the send idempotent per
recipient (Section~\ref{sec:enhancing-kafka}).
\end{pitfall}

\section{Pros and cons, candidly}

\begin{center}
\begin{tabular}{@{}p{0.28\linewidth}p{0.32\linewidth}p{0.32\linewidth}@{}}
\toprule
 & The project today & Production baseline \\
\midrule
Producer durability & \code{acks=all}, 1 replica & 3 replicas, \code{min.insync}=2 \\
Producer duplicates & none (idempotent) & none (idempotent) \\
Consumer duplicates & possible, unhandled & possible, deduplicated \\
Failed record & log and skip after 10 & retry topics, then DLT \\
Ordering & per key, 1 partition & per key, N partitions \\
Slow work & on the poll thread & fanned out, or async \\
Observability & none on lag & lag alerts, DLT alerts \\
\bottomrule
\end{tabular}
\end{center}

The project's honest weak points: silent loss after ten retries; e-mail
fan-out on the poll thread; no dedup where the side effect is not
naturally idempotent (mail); no lag metric anyone looks at; and
single-partition topics, which make ordering trivial and parallelism
impossible.

\section{Enhancing the resolution}
\label{sec:enhancing-kafka}

\subsection{An idempotent consumer with a record of what it has seen}

\code{UserActivatedListener} is already idempotent by natural key --- a
second delivery finds the student and returns. Where there is no
natural key, record the event id. The outbox of Chapter~33 puts a
unique \code{eventId} header on every message; the consumer keeps a
table of ids it has finished, and writes that row \emph{in the same
transaction} as the business change. Flyway
\code{V3\_\_create\_processed\_events.sql} in course-service:

\begin{lstlisting}[language=SQL]
CREATE TABLE processed_events (
    event_id     UUID PRIMARY KEY,
    consumer     VARCHAR(64) NOT NULL,     -- listener name
    processed_at TIMESTAMPTZ NOT NULL DEFAULT now()
);
-- retention: rows older than the topic's retention are useless
CREATE INDEX idx_processed_events_at ON processed_events (processed_at);
\end{lstlisting}

The listener, rewritten to read the header and to claim the id first:

\begin{lstlisting}[language=Java]
@KafkaListener(topics = "ts.user.activated", groupId = "course-group")
@Transactional                                              (*@\co{1}@*)
public void onUserActivated(
        UserActivatedEvent event,
        @Header("eventId") String eventId) {

    int claimed = jdbc.update("""
        INSERT INTO processed_events (event_id, consumer)
        VALUES (?::uuid, 'UserActivatedListener')
        ON CONFLICT (event_id) DO NOTHING""", eventId);      (*@\co{2}@*)
    if (claimed == 0) {
        log.info("duplicate event {} ignored", eventId);
        return;                                             (*@\co{3}@*)
    }
    if (!ROLE_STUDENT.equalsIgnoreCase(event.role())) {
        return;
    }
    Long userId = Long.parseLong(event.userId());
    Student student = studentRepository.findByUserId(userId)
            .orElseGet(() -> Student.builder().userId(userId).build());
    student.setFirstName(event.firstName());               (*@\co{4}@*)
    student.setLastName(event.lastName());
    student.setEmail(event.email());
    studentRepository.save(student);
}
\end{lstlisting}

\begin{description}
  \item[\co{1}] One database transaction around claim and change. If
  \code{save()} fails, the claim rolls back too, and the retry can claim
  again. If the JVM dies after commit but before the Kafka offset
  commit, the redelivery hits \co{3}.
  \item[\co{2}] \code{ON CONFLICT DO NOTHING} is the whole dedup
  algorithm: an atomic insert-if-absent. No read-then-write race between
  two consumer instances that got the same record after a rebalance.
  \item[\co{3}] Return normally, not by throwing: the offset must
  advance past a duplicate.
  \item[\co{4}] While here, upsert instead of insert-if-missing, so the
  same listener also serves the \code{ts.user.updated} event Chapter~33
  adds.
\end{description}

Notification-service has no database, and mail is the one side effect
you cannot un-send. Options, honest about cost: accept an occasional
duplicate mail (what the project implicitly does); or give the service a
tiny store --- Redis \code{SET eventId 1 NX EX 604800} before sending,
which is the same claim in one command; or make the mail itself
idempotent by including the event id in a \code{Message-ID} header, so
a mailbox de-duplicates what the consumer could not.

\subsection{Retry topics and a dead-letter topic}

Replace ``retry ten times in a millisecond, then forget'' with retries
that wait, and a destination for what still fails. Spring Kafka's
\code{@RetryableTopic} does this with topics instead of threads, so a
stuck record does not block the partition behind it:

\begin{lstlisting}[language=Java]
@RetryableTopic(
        attempts = "4",                                     (*@\co{1}@*)
        backoff = @Backoff(delay = 2_000, multiplier = 3),  (*@\co{2}@*)
        include = { MailSendException.class,
                    MailAuthenticationException.class },    (*@\co{3}@*)
        dltStrategy = DltStrategy.FAIL_ON_ERROR)
@KafkaListener(topics = "ts.user.registered",
               groupId = "notification-group")
public void onUserRegistered(UserRegisteredEvent event) {
    emailService.sendActivationEmail(
            event.email(), event.firstName(), event.activationCode());
}

@DltHandler                                                 (*@\co{4}@*)
public void onDeadLetter(UserRegisteredEvent event,
        @Header(KafkaHeaders.EXCEPTION_MESSAGE) String reason,
        @Header(KafkaHeaders.ORIGINAL_OFFSET) long offset) {
    log.error("DLT ts.user.registered offset={} user={} reason={}",
            offset, event.userId(), reason);
    meter.counter("kafka.dlt", "topic", "ts.user.registered")
         .increment();                                      (*@\co{5}@*)
}
\end{lstlisting}

\begin{description}
  \item[\co{1}] One try on the main topic, three on retry topics.
  \item[\co{2}] Delays of 2, 6, and 18 seconds --- about half a minute
  for MailDev or an SMTP relay to come back. Spring creates the topics
  \code{ts.user.registered-retry-2000}, \code{-retry-6000},
  \code{-retry-18000} and \code{ts.user.registered-dlt}; under Strimzi
  (Chapter~25) declare them as \code{KafkaTopic} objects or leave
  auto-creation on.
  \item[\co{3}] Retry only what a wait can fix. A deserialization
  failure or a \code{NullPointerException} goes straight to the DLT;
  retrying it three times is theatre.
  \item[\co{4}] The dead-letter handler is a normal listener on the
  \code{-dlt} topic. It \emph{must not} throw.
  \item[\co{5}] A counter, so that ``DLT non-empty'' is an alert
  (Chapter~27's observability) instead of a discovery.
\end{description}

What to do with a dead letter is an operational decision, not code: an
alert fires; someone reads the \code{EXCEPTION\_MESSAGE} header; the fix
ships; and the records are \emph{replayed} --- copied from the DLT back
to the main topic with a small admin tool or Kafka UI --- where the
idempotent consumer guarantees the ones that did succeed are ignored.

\begin{handson}[see the retry ladder and the dead letter]
Break SMTP, register a user, and read the topics:
\begin{lstlisting}[language=bash]
$ docker compose stop maildev
$ curl -s -X POST localhost:8080/api/auth/register -H \
  'Content-Type: application/json' -d '{"email":"dlt@test.io", \
  "password":"Passw0rd!","firstName":"D","lastName":"L", \
  "activationMethod":"EMAIL"}'
$ docker compose exec kafka /opt/kafka/bin/kafka-topics.sh \
  --bootstrap-server localhost:19092 --list | grep registered
ts.user.registered
ts.user.registered-dlt
ts.user.registered-retry-18000
ts.user.registered-retry-2000
ts.user.registered-retry-6000
\end{lstlisting}
About 26 seconds after the request, notification-service logs
\code{DLT ts.user.registered offset=0 user=... reason=Mail server
connection failed}. Start MailDev again, copy the record from the DLT to
\code{ts.user.registered} in Kafka UI (\code{localhost:9090}), and the
activation mail arrives.
\end{handson}

\subsection{Exactly-once, and what it does not cover}

Kafka's transactions give exactly-once for one specific shape:
consume from topic A, produce to topic B, commit both offsets and
output atomically (\code{isolation.level=read\_committed} on the reader).
Stream processors need it; the project has no such stage. Transactions
do \emph{not} span Kafka and Postgres, or Kafka and SMTP --- the moment
the side effect leaves the broker, you are back to at-least-once plus
idempotency. An interviewer who asks ``can you get exactly-once?''
wants exactly that distinction: yes inside Kafka, no across systems,
and the outbox plus a dedup table is how you get the effect across
systems anyway.

\section{What breaks at 3\,a.m.}

\begin{itemize}
  \item \textbf{notification-service down for four hours.} Nothing is
  lost: records accumulate on the topic, retention is 168 hours in
  Compose and Strimzi's default seven days. Consumer lag ---
  the gap between the partition's last offset and the group's committed
  offset --- grows to a few thousand. On restart the service drains them
  at SMTP speed; every activation mail is late, none is missing. What
  \emph{would} lose them: an outage longer than retention, or a
  redeploy under a \emph{new} group id, which starts from
  \code{earliest} and resends every mail in seven days of history.
  \item \textbf{The broker is down.} Producers block for
  \code{max.block.ms} on their first metadata fetch, then throw --- the
  sixty-second registration failure of Chapter~33. Producers that were
  already connected buffer up to \code{buffer.memory} (32\,MB) and retry
  within \code{delivery.timeout.ms}; after two minutes the future fails
  and, since nobody reads it, the event is gone. Consumers simply poll
  nothing. The outbox converts all of this into a growing
  \code{attempts} column.
  \item \textbf{A replay from yesterday.} An operator resets
  \code{course-group} to a timestamp with
  \code{kafka-consumer-groups.sh --reset-offsets --to-datetime}. Every
  event since then is redelivered; with the \code{processed\_events}
  claim each one costs one failed insert. Without it, mail goes out
  twice and student rows are attempted twice (saved only by the
  \code{UNIQUE(user\_id)} constraint).
  \item \textbf{A hot key.} Every event for one very large course lands
  on one partition; adding partitions does not help because the key
  chooses the partition. Re-key --- by recipient instead of by
  assignment --- or accept that one key's throughput is one consumer's
  throughput.
\end{itemize}

\section{Outside the project}

Every managed Kafka --- Confluent Cloud, AWS MSK, Aiven --- ships
\code{acks=all}, three replicas, and \code{min.insync.replicas=2} as the
starting point, and sells lag dashboards as the first thing you look at.
Retry topics and DLTs are so standard that Confluent's Parallel
Consumer, Uber's consumer proxy, and Spring's \code{@RetryableTopic}
all implement the same ladder; the differences are in how they preserve
per-key ordering across retries (Spring does not, by default --- a
retried record can be processed after a later one for the same key,
which is fine for mail and not for ledger entries). RabbitMQ and SQS
give you per-message acknowledgement and a dead-letter queue as a broker
feature rather than a library pattern; what they cannot give you is
replay, which is the reason a log was chosen in the first place.

\section{Questions from real interviews}

\subsection*{Q: How do you increase the number of messages Kafka can handle per unit of time?}

I start with the consumer, because that is where the project's limit
is. Partitions are the unit of parallelism: a group cannot have more
active consumers than the topic has partitions, and every topic in
\code{kafka-topics.yaml} has one. So the first change is partition
count --- 6 or 12, chosen with headroom because partitions can be added
but never removed, and adding them later changes the key-to-partition
hash, which breaks per-key ordering for in-flight data. Then match
consumers to partitions: \code{spring.kafka.listener.concurrency=3} per
pod times two replicas gives six threads over six partitions. Next,
get slow work off the poll thread: one \code{AssignmentCreated} that
sends 300 e-mails synchronously is a five-minute batch that trips
\code{max.poll.interval.ms}; publish one record per recipient instead,
so throughput scales with partitions. On the producer, batching:
\code{linger.ms=5--20} and \code{batch.size=64--256\,KB} let the client
send hundreds of records per request, and
\code{compression.type=lz4} or \code{zstd} cuts bytes on the wire by
three to five times for JSON. \code{acks=1} roughly doubles producer
throughput at the cost of losing a write when the leader dies before
replicating --- I keep \code{acks=all} for anything I cannot
regenerate. On the consumer, \code{max.poll.records} and
\code{fetch.min.bytes} trade latency for fewer round-trips, bounded by
what fits in \code{max.poll.interval.ms}. And I watch for a hot key:
if one partition carries 80\,\% of the load because of a bad key
choice, none of the above helps and the fix is a different key.

\subsection*{Q: Consumer lag is growing. How do you diagnose it?}

First I establish which of three things it is: the consumer is slow,
the consumer is dead, or the producer got faster. \code{kafka-consumer-
groups.sh --describe --group notification-group} shows lag per
partition and which member owns each; lag on \emph{all} partitions with
members present means slow processing; lag on \emph{some} with missing
members means an instance is down or stuck in rebalance; lag that
appeared at a clear time with steady consumption means a producer
burst. For slow processing I look at the listener's per-record time
--- in this project SMTP latency --- and at rebalance churn in the logs
(\code{Member ... has left}), which means \code{max.poll.interval.ms} is
being hit and the group is thrashing. Fixes follow the diagnosis: more
partitions and concurrency, move the slow call off the thread, or,
if the burst is legitimate, let it drain and alert only when lag is
older than the SLO in seconds, not larger than a number of records. The
metric to graph is lag in \emph{time}, not count.

\subsection*{Q: A message was processed twice. Why, and how do you make that safe?}

Because Kafka commits offsets after processing, and the process died,
or was evicted for being slow, between the two. That is at-least-once
working as designed. Making it safe means every consumer either has a
natural idempotency key --- \code{UserActivatedListener}'s
\code{findByUserId} --- or claims a unique event id in the same
transaction as its side effect: \code{INSERT ... ON CONFLICT DO
NOTHING} into \code{processed\_events}, return if zero rows. The id has
to come from the producer, minted once and carried through retries,
which is what the outbox row gives you. For side effects with no
transaction --- e-mail --- I put the id in the message's
\code{Message-ID} or claim it in Redis with \code{SET NX}. What I do
not do is switch to \code{AckMode.RECORD} and call it fixed; that
narrows the window from a batch to one record and leaves the bug in
place.

\subsection*{Q: How do you guarantee ordering for one user?}

Key every event about that user by the user id, so the hash puts them
on one partition, and keep one consumer thread per partition. That is
the whole guarantee, and it has three ways to break: an idempotent
producer disabled with \code{max.in.flight} above one, which lets a
retried batch land after a later one; a partition count change, which
re-hashes keys; and retry topics, which re-inject a failed record
behind its successors. For the project, \code{userId} keys on
\code{ts.user.*} already do this. If the business needs a stronger
statement --- ``the update is applied after the create, always'' --- I
add a version number to the event and have the consumer reject anything
older than what it has, which is what turns ordering from a broker
property into a consumer invariant.

\subsection*{Q: A broker dies. What does the producer do?}

With three brokers and \code{replication.factor=3}, the controller
elects a new leader from the in-sync replicas for each affected
partition within seconds; the producer gets
\code{NotLeaderOrFollower}, refreshes metadata, and retries ---
invisible to the application because \code{retries} is unbounded and
\code{delivery.timeout.ms} is two minutes. If the dead broker took the
ISR below \code{min.insync.replicas}, the producer gets
\code{NotEnoughReplicas} and keeps retrying; with \code{acks=all} that
is the correct refusal to write a single copy. On the project's
single-broker Compose stack there is no failover: sends buffer, the
future fails after two minutes, and the outbox row is what saves the
event. The one thing I check in a review is whether anyone reads the
future --- in this project nobody does, so a producer-side failure is
silent.

\subsection*{Q: How do you replay a topic from yesterday?}

Stop the consumer group, reset its offsets, start it again:
\code{kafka-consumer-groups.sh --reset-offsets --group course-group
--topic ts.user.activated --to-datetime 2026-09-06T00:00:00.000
--execute}. The group must be inactive or the command refuses. Before
running it I answer two questions: is the consumer idempotent for
everything since that time --- because every record comes back --- and
is the data still there, which depends on retention (\code{168h} here)
and on whether the topic is compacted. For a single record I replay
from the DLT instead of resetting the whole group. And I never replay
into a group that sends e-mail without first confirming the dedup path,
because ``we resent 40\,000 activation mails'' is the incident that
follows a confident reset.

\begin{quiz}
\begin{enumerate}
  \item Name the producer settings that make the project's producers
  idempotent and durable, and explain why \code{acks=all} on the Compose
  stack is weaker than it sounds.
  \item Draw (in words) the timeline in which \code{notification-group}
  sends the same activation mail twice without any crash, and name the
  setting whose value defines the trigger.
  \item Why must the \code{processed\_events} insert and the business
  write share one transaction? Describe the interleaving that goes
  wrong if the insert commits first on its own.
  \item In the \code{@RetryableTopic} configuration, explain why
  \code{include} lists only mail exceptions and what happens to a
  \code{NullPointerException}.
  \item Course-service declares its own \code{ProducerFactory} bean.
  State the consequence for the \code{spring.kafka.producer} block in
  its YAML, and the one-line fix.
  \item Notification-service has been down for nine days. State what
  happens on restart, which setting decides it, and what would be
  different if it restarted with \code{group-id: notification-v2}.
\end{enumerate}
\end{quiz}
